% 全集实验的方法段落草稿, 尚未并入主稿.
% 本文件不含性能结果. 全部任务完成并核验后, 再结合结果表精简并入 exp-evaluation.tex.
% 最终并入前重新核对全部执行节点, 并使用最终实验归档提交替换主稿的工件版本指针.
% 引用均已存在于当前主稿; 不修改 references.bib 或摘要与关键词.

\paragraph{公开基准的范围.}
公开基准采用三个研究工件的完整可获取输入集合: Trentin 与 Sebastiani~\cite{DBLP:journals/jar/TrentinS21} 的浮点优化基准, Nadel 与 Ryvchin~\cite{DBLP:conf/tacas/NadelR16} 的位向量优化基准, 以及 He 等人~\cite{DBLP:conf/fm/HeLZC24} 发布的 MaxSMT(LIA) 基准. 固定数据版本后逐文件建立清单, 保留各来源组和目标方向, 不作实例抽样. 数据规模按源文件和实际优化任务分别统计.

浮点工件的原始 \texttt{bench} 目录包含 $17\,314$ 个源文件, 分为 griggio、heizmann、ramalho、schanda、wintersteiger 五组. 发布的 \texttt{o300} 集合含 $1120$ 项, 是经过规模限制的子集. 对完整源集合沿用原工件的目标构造, 保留标记为不可满足的源文件, 得到 $39\,638$ 个任务; 补齐原生成器未识别的目标后增加 $16$ 项. \texttt{o300} 中另有一个内容不同的输入, 一并保留, 合计 $39\,655$ 项. 其余 $1119$ 个发布输入与生成结果逐字节相同, 记录其对应关系而不重复计数. 补充项中有两个来自空公式源, 以 binary32 常量正零分别构造最小化与最大化任务; 二者标为退化检查, 不用于说明浮点运算的求解特点. 该任务集合覆盖全部原始源文件和全部发布输入.

位向量工件包含 $254$ 个原生 BV 输入: 五个布局规模组各 $40$ 项, 另有 \texttt{hd} 组 $54$ 项. 全部输入保留原硬约束, 按无符号次序优化原目标. 该公开集合包括布局问题的构造实例, 不等同于原文所述的非公开工业实例全集. 原包另附的 $254$ 个整数文件全部作为多目标语义诊断, 保留原约束与目标顺序并显式采用字典序. 发布包中至少一个整数伴随文件的目标可上无界增长, 而对应 BV 目标只有有限取值; 因此不将两组文件作为等价编码作求解速度比较.

MaxSMT 工件的标准 SMT-LIB2 输入共 $2264$ 项, 由 $283$ 个基础问题、四种软约束比例及有权/无权两种设置组合而成. 全部标准输入进入清单, 另一输入格式中的副本不重复计数. 该公开包包含 MaxSMT(LIA) 输入, 其范围不等同于原文报告的全部实验集合. 全量解析中 $2245$ 项通过, $19$ 项在发布包中已有语法错误或文件截断. 后者保留原文件和求解器的实际响应, 单列为无效输入; 完整发布规模与语法有效规模分别报告.

\paragraph{公开基准的求解配置.}
FP 比较三种配置: OptiMathSAT 直接优化原生 FP 目标; Z3 保留 FP 约束, 以 BV 序键编码目标; Z3 进一步使用 \texttt{fpa2bv} 将 FP 运算展开为 BV 约束. 三种配置均排除 NaN 目标. 对非 NaN 目标的 $w$ 位 IEEE 表示 $b$, 序键为
\[
k(b)=\begin{cases}
  \mathop{\mathrm{bvnot}}b,& b_{w-1}=1,\\
  b\mathbin{\mathrm{bvxor}}2^{w-1},& b_{w-1}=0.
\end{cases}
\]
序键保持非 NaN 数值顺序, 并将有符号零细化为 $-0<+0$. 比较原生数值目标与序键目标的最优值时, 将两个零视为数值相等.

原生 BV 比较 OptiMathSAT 与 Z3 的无符号位向量优化. MaxSMT 比较 OptiMathSAT 的 MaxRes、OMT 两种配置与 Z3 的 MaxRes、WMax 两种配置, 各配置使用相同的硬约束、软约束及权重. 整数伴随诊断使用 OptiMathSAT 和 Z3, 对两者均显式设置字典序, 不由某个工具的默认多目标模式决定问题语义.

\paragraph{运行环境、预算与统计.}
公开基准在 Linux CPU 机群上运行, 与前述 WSL 参数化实验分别统计. 执行节点采用 Ubuntu 22.04、24.04 或 26.04, 处理器包括 Xeon Platinum 8153、Xeon Gold 6248、Xeon Gold 5215 和 EPYC 7702. 同一输入的各配置在同一节点依次执行, 执行次序由输入标识预先确定, 与运行结果无关. Z3 与 OptiMathSAT 分别使用 4.15.4、1.7.4, 各节点使用相同的求解器二进制版本. 逐次运行的处理器、系统、Python 版本与二进制校验值保存在实验工件中.

每次求解采用单线程、$600$\,s 墙钟预算和 $8$\,GiB 求解进程地址空间上限; 包含输入处理程序的整个作业上限为 $16$\,GiB. 求解计时包含输入解析、目标转换与进程启动, 独立复核时间另计. 对 FP、BV 和 MaxSMT 返回的候选目标值, 在原始硬约束下分别检查可达性与严格改进, 每条查询限时 $60$\,s. 只有前者返回 SAT、后者返回 UNSAT, 才记为独立复核通过. 求解器报告最优的数量与复核通过的数量分别统计, 复核超时不记为通过. 整数伴随文件仅报告状态与界信息, 不将 SAT 或打印出的模型解释为已证明最优.

完整清单中保留全部超时、内存不足、未定状态和无效输入. 性能时间统计以语法有效输入为分母; 求解器报告最优或不可满足时记为完成判定, 其余有效输入按 $+\infty$ 参与中位数计算. PAR-2 为有效输入上的平均耗时, 未完成判定的输入按两倍时限计入. 运行失败的分类、原始输出、最优值一致性检查和独立复核结果均与输入清单逐项对应.
